\section{GStreamer}

GStreamer es un \textit{framework} multimedia de código abierto, multiplataforma y
modular, publicado bajo licencia LGPL. Su arquitectura se basa en el patrón
\textit{pipeline}: los elementos de procesamiento (fuentes, filtros, transformaciones,
sumideros) se interconectan mediante \textit{pads} tipados, formando grafos dirigidos
acíclicos por los que fluyen los búferes de datos multimedia. Esta arquitectura permite
construir cadenas de procesamiento arbitrariamente complejas combinando elementos
existentes sin necesidad de escribir código de bajo nivel.

Los \textit{caps} (\textit{capabilities}) describen el tipo de datos que cada pad
acepta o produce (formato de píxel, frecuencia de muestreo, \textit{framerate}, etc.),
y GStreamer los negocia automáticamente en tiempo de construcción del pipeline,
garantizando la compatibilidad entre elementos adyacentes.

GStreamer es la base sobre la que se construye Voctomix: actúa como motor de bajo
nivel para toda la manipulación de señal: decodificación, composición, conversión de
formato y transmisión de red.

El elemento central en Voctomix 2.0 es \texttt{compositor}, el mezclador de vídeo
multipista de GStreamer 1.x, que reemplaza al antiguo \texttt{videomixer}. Este
elemento acepta $N$ pads de entrada, cada uno con su propia transformación geométrica
(posición, escala, transparencia), y produce un único flujo de salida. Las propiedades
de cada pad se modifican en tiempo real mediante mensajes GObject, lo que permite
implementar transiciones suaves entre composites sin interrumpir el pipeline ni generar
fotogramas corruptos.

Para la sincronización entre el núcleo y múltiples clientes remotos, Voctomix utiliza
\texttt{GstNetTimeProvider} (puerto UDP 9998), un reloj de red que garantiza que todos
los nodos del sistema operen sobre la misma referencia temporal, condición indispensable
para evitar artefactos de sincronización audio-vídeo en escenarios multi-equipo.

% --- Figura sugerida ---
% \begin{figure}[H]
%   \centering
%   \includegraphics[width=0.85\textwidth]{figuras/pipeline_gstreamer.png}
%   \caption{Pipeline GStreamer simplificado de voctocore (fuentes, compositor y salidas).}
%   \label{fig:pipeline_gstreamer}
% \end{figure}

% --- Tabla sugerida ---
% \begin{table}[H]
%   \centering
%   \caption{Elementos GStreamer utilizados en Voctomix 2.0 y su función.}
%   \label{tab:elementos_gstreamer}
%   \begin{tabular}{|l|l|}
%     \hline
%     \textbf{Elemento} & \textbf{Función} \\
%     \hline
%     \texttt{compositor}        & Composición multipista de vídeo \\
%     \texttt{matroskamux}       & Encapsulado MKV para transporte TCP \\
%     \texttt{tcpclientsrc}      & Recepción de fuente por socket TCP \\
%     \texttt{tcpserversink}     & Publicación de salida por socket TCP \\
%     \texttt{gdkpixbufoverlay}  & Superposición de imagen PNG (overlays) \\
%     \texttt{volume}            & Control de ganancia de audio \\
%     \texttt{audiomixer}        & Mezcla multipista de audio \\
%     \hline
%   \end{tabular}
% \end{table}
